Un mapa topográfico de una montaña muestra curvas de nivel equiespaciadas en altura ($\Delta z = 100$ m) a distancias horizontales que varían de un punto a otro. En la zona $A$ las curvas están muy juntas; en la zona $B$ están muy separadas.
\begin{enumerate}
    \item ¿En cuál zona es más pronunciada la pendiente del terreno? Justificar usando la interpretación de las curvas de nivel.
    \item Si $z = f(x,y)$ modela la altitud, y en la zona $A$ una curva de nivel $c=500$ m y la siguiente $c=600$ m están separadas apenas $50$ m horizontalmente, mientras que en la zona $B$ están separadas $500$ m, estimar la ``pendiente media'' en cada zona como $\Delta z / \Delta d$.
    \item Proponer una función $f(x,y)$ sencilla cuyas curvas de nivel tengan la forma general de un volcán (alto en el centro, nivel descendente hacia afuera).
\end{enumerate}